\documentclass[a4paper]{article}
% Español (México) con XeLaTeX: fontspec para fuentes Unicode, polyglossia
% para la división silábica y los nombres del documento en la variante mexicana.
\usepackage{fontspec}
\usepackage{polyglossia}
\setdefaultlanguage[variant=mexican]{spanish}
% Los nombres chinos van como «pinyin (original)»: xeCJK da a los caracteres
% Han su propia fuente; sin ella XeLaTeX los omite con un aviso.
% FandolSong no cubre algunos caracteres de nombres (U+73FA); WenQuanYi Zen Hei sí.
\usepackage[AutoFallBack=true]{xeCJK}
\setCJKmainfont{FandolSong-Regular.otf}
\setCJKfallbackfamilyfont{\CJKrmdefault}{WenQuanYi Zen Hei}
% polyglossia no traduce los nombres de listings.
\AtBeginDocument{\providecommand{\lstlistingname}{}\renewcommand{\lstlistingname}{Listado}%
  \providecommand{\lstlistlistingname}{}\renewcommand{\lstlistlistingname}{Índice de listados}}
\usepackage{amsmath, amssymb}
\usepackage{graphicx}
\usepackage[margin=2.35cm]{geometry}
\usepackage[most]{tcolorbox}
\usepackage{booktabs}
\usepackage{tabularx}
\usepackage{longtable}
\usepackage{array}
\usepackage{float}
\usepackage{xcolor}
\usepackage{hyperref}
\usepackage{fancyhdr}
\setCJKmainfont[AutoFakeBold=2.4]{AR PL UMing CN}
\setCJKsansfont[AutoFakeBold=2.4]{Droid Sans Fallback}
\setCJKmonofont[AutoFakeBold=2.4]{Droid Sans Fallback}
\hypersetup{colorlinks=true, linkcolor=blue!55!black, urlcolor=blue!60!black}
\graphicspath{{./}{figures/}}
\setlength{\parskip}{0.55em}
\setlength{\parindent}{2em}
\setlength{\emergencystretch}{3em}
\renewcommand{\arraystretch}{1.18}
\newtcolorbox{knowledgebox}[1]{enhanced, colback=blue!5!white, colframe=blue!70!black, colbacktitle=blue!70!black, coltitle=white, fonttitle=\bfseries, title={#1}, sharp corners, breakable}
\newtcolorbox{importantbox}[1]{enhanced, colback=yellow!10!white, colframe=yellow!75!black, colbacktitle=yellow!75!black, coltitle=black, fonttitle=\bfseries, title={#1}, sharp corners, breakable}
\newtcolorbox{warningbox}[1]{enhanced, colback=red!5!white, colframe=red!70!black, colbacktitle=red!70!black, coltitle=white, fonttitle=\bfseries, title={#1}, sharp corners, breakable}
\pagestyle{fancy}
\fancyhf{}
\fancyfoot[C]{\thepage}
\fancyfoot[R]{\footnotesize\color{gray}\href{https://github.com/hqhq1025/ai-course-notes}{AI Course Notes}}
\renewcommand{\headrulewidth}{0pt}
\renewcommand{\footrulewidth}{0pt}
\newcommand{\videourl}{https://www.youtube.com/watch?v=6yExfoTuSWw}
\newcommand{\repourl}{https://github.com/hqhq1025/ai-course-notes}

\begin{document}
\begin{titlepage}
\centering
{\Large Notas didácticas de revisión técnica\par}
\vspace{1.0cm}
{\huge\bfseries Informe trimestral de los grandes modelos: diferenciación, convergencia, suites integrales y la experiencia L4}\par
\vspace{0.5cm}
{\Large Zhang Xiaojun conversa con Guangmi (广密)\par}
\vspace{0.35cm}
{\large Elaborado a partir del video público de Zhang Xiaojun Podcast, la transcripción generada en GPU, la estructura de capítulos y diagramas conceptuales propios\par}
\vspace{0.3cm}
{\large 2025-08-18\par}
\vspace{0.85cm}
\includegraphics[width=0.88\textwidth,height=0.42\textheight,keepaspectratio]{cover.jpg}\par
\vfill
\begin{tcolorbox}[width=0.92\textwidth, colback=black!2!white, colframe=black!55, sharp corners]
\textbf{Título del video}: 112. Conversación con Guangmi sobre el informe trimestral de los grandes modelos: diferenciación y convergencia, suites integrales e integración vertical, la experiencia L4 y la ventana de minería\par
\textbf{Canal del video}: Zhang Xiaojun Podcast\par
\textbf{Duración del video}: 01:09:11\par
\textbf{Enlace del video}: \href{\videourl}{\nolinkurl{\videourl}}\par
\textbf{Notas sobre el material}: YouTube no tiene subtítulos disponibles; el texto principal se elaboró a partir de la transcripción con faster-whisper large-v3 en CUDA, los capítulos del video, la descripción del video, la portada y figuras didácticas propias. El video de portada fija no se reutiliza en el cuerpo del texto.\par
\textbf{Página del proyecto}: \href{\repourl}{\nolinkurl{\repourl}}\par
\end{tcolorbox}
\end{titlepage}

\tableofcontents
\newpage

\section{Introducción: las dos palabras clave de esta temporada}

Esta sección establece primero el hilo conductor de todo el informe trimestral. Las ediciones anteriores del «Informe trimestral global de modelos grandes» se centraban más en el límite superior de la inteligencia, el avance de los modelos y la línea de la AGI; esta edición empieza a orientarse con claridad hacia el producto, los puntos de entrada y la comercialización. Las dos palabras clave que propone Guang Mi (广密) son «diferenciación» y «producto». La diferenciación se refiere a que las empresas de modelos empiezan a separarse en distintos rangos de capacidad y distintas rutas de producto; el producto se refiere a que la inteligencia del modelo tiene que convertirse en experiencia de usuario, tráfico, presencia de marca en la mente del usuario y un ciclo comercial cerrado.

La tensión central de esta edición es la siguiente: por un lado, los tres primeros modelos base siguen avanzando con rapidez y absorben continuamente las oportunidades de los emprendedores de IA; por otro lado, la ventana de producto del Magic Moment todavía existe, aunque su vida útil es cada vez más corta. Los emprendedores enfrentan una situación a la vez emocionante y desesperante: la tecnología se vuelve más potente cada mes, pero los gigantes también entran al mercado cada vez más rápido.

\begin{figure}[H]
\centering
\includegraphics[width=0.96\textwidth]{startup-between-giants.png}
\caption{Los emprendedores, atrapados entre los gigantes: la tecnología avanza cada mes, pero las empresas de modelos también siguen absorbiendo oportunidades. Diagrama conceptual propio, elaborado a partir de la conversación de 00:00:05--00:03:54.}
\end{figure}

\begin{knowledgebox}{Lectura de la figura: por qué es emocionante y desesperante a la vez}
Lo emocionante proviene del avance tecnológico; lo desesperante, de que las principales empresas de modelos desarrollan al mismo tiempo los modelos y los productos. Los emprendedores tienen que encontrar la Magic Window antes de que los gigantes conviertan esas capacidades en producto, o bien elegir escenarios verticales que a los gigantes les resulte difícil cubrir con rapidez.
\end{knowledgebox}

\begin{importantbox}{Tesis central de esta edición}
La industria de los modelos grandes pasa de una «competencia por el límite superior de la inteligencia» a una competencia compuesta de «inteligencia + producto + punto de entrada + comercialización». La capacidad del modelo sigue siendo importante, pero las barreras de producto, el punto de entrada predeterminado, la presencia de marca en la mente del usuario, la integración vertical y la experiencia L4 empiezan a decidir quién logra convertir la inteligencia en valor real.
\end{importantbox}

\subsection{Resumen de la sección}

EP112 es el episodio en que el informe trimestral de modelos pasa de la «inteligencia del modelo» a la «estrategia de producto». No se limita a enumerar la actividad de las empresas: responde a la pregunta de dónde pueden apostar todavía los emprendedores y los inversionistas una vez que las principales empresas de modelos empiezan a convertir sus capacidades en producto.
\section{Los modelos se diferencian: generalistas, Coding+Agentic e interacción multimodal}

La introducción planteó dos palabras clave, «diferenciación» y «producto»; esta sección desarrolla la primera. La cuestión central es que las principales empresas de modelos ya no compiten solo sobre la misma curva de benchmark, sino que empiezan a separarse en su pila de capacidades y en su trayectoria de producto. Este trimestre, las empresas de modelos de Silicon Valley comenzaron a diferenciarse de manera notable. Google Gemini y OpenAI siguen apostando por las capacidades generales; Anthropic se diferenció hacia las capacidades de modelo en Coding y Agentic; Thinking Machines Lab se diferenció hacia lo multimodal y la interacción de nueva generación; Grok todavía busca su nicho en el ecosistema; y se considera que Meta tiene un ADN débil para la creación original de 0 a 1. Las empresas más avanzadas se parecen a una competencia de F1: avanzan a una velocidad altísima y sus rutas también empiezan a divergir.

\begin{figure}[H]
\centering
\includegraphics[width=0.96\textwidth]{model-company-divergence.png}
\caption{Las empresas de modelos empiezan a diferenciarse: Gemini/OpenAI generalistas, Anthropic Coding+Agentic, Thinking Machines interacción multimodal. Diagrama conceptual propio, elaborado a partir de la conversación entre 00:03:54 y 00:21:37.}
\end{figure}

\begin{knowledgebox}{Lectura de la figura: diferenciarse no es quedarse atrás, sino elegir un carril}
Los modelos generalistas deben mantener el liderazgo en todas las capacidades; Coding+Agentic apuesta por flujos de trabajo de alto valor; lo multimodal y la interacción apuestan por el punto de entrada de la próxima generación; las empresas centradas en su ecosistema necesitan encontrar cómo integrarse con su propia plataforma. La diferenciación muestra que la industria pasó de perseguir únicamente benchmarks a tomar decisiones estratégicas.
\end{knowledgebox}

\begin{knowledgebox}{Términos clave: mapa de rutas de las empresas}
\begin{tabularx}{\textwidth}{p{0.22\textwidth}p{0.34\textwidth}X}
\toprule
Empresa/ruta & Posicionamiento en este episodio & Riesgo principal \\
\midrule
Gemini/OpenAI & Modelos generalistas y punto de entrada con todas las capacidades & Costo alto y líneas de producto complejas. \\
Anthropic & Capacidades de Coding + Agentic & Debe convertir su ventaja técnica en un punto de entrada de producto más amplio. \\
Thinking Machines & Multimodalidad nativa e interacción de nueva generación & Dirección novedosa, pero con un camino comercial aún indefinido. \\
Grok/xAI & Todavía busca su nicho en el ecosistema & La sinergia con la plataforma y el posicionamiento de marca aún deben demostrarse. \\
Meta & Ecosistemas de código abierto y redes sociales sólidos, pero se cuestiona su capacidad de creación original de 0 a 1 & Inestabilidad entre su ruta de modelos y su punto de entrada de producto. \\
\bottomrule
\end{tabularx}
\end{knowledgebox}

\begin{knowledgebox}{La pila de capacidades detrás de la diferenciación}
\begin{tabularx}{\textwidth}{p{0.22\textwidth}p{0.34\textwidth}X}
\toprule
Pila de capacidades & Empresas o direcciones representativas & Activos principales \\
\midrule
Inteligencia general & Gemini, OpenAI & Multimodalidad, razonamiento, búsqueda, puntos de entrada de producto y capacidad de cómputo. \\
Coding/Agentic & Anthropic, Claude Code & Entornos de código, tareas de largo horizonte y posicionamiento entre desarrolladores empresariales. \\
Interacción multimodal & Thinking Machines, Gemini & Multimodalidad nativa, interface de nueva generación y hardware como punto de entrada. \\
Modelos de ecosistema & Grok, Meta & Plataformas sociales, ecosistema de código abierto y canales de distribución. \\
\bottomrule
\end{tabularx}
\end{knowledgebox}

\begin{knowledgebox}{Cuadro comparativo de estrategias de las empresas de modelos}
\begin{tabularx}{\textwidth}{p{0.18\textwidth}p{0.30\textwidth}p{0.40\textwidth}}
\toprule
Empresa/ruta & Fortalezas actuales & Lo que aún debe demostrar \\
\midrule
OpenAI & Punto de entrada de ChatGPT, capacidades generales, conversión en producto & Si la publicidad y la comercialización pueden convivir sin dañar la experiencia del usuario. \\
Google & TPU, búsqueda, publicidad, Workspace, Gemini & Si puede recombinar sus activos para convertirse en el punto de entrada predeterminado de la IA. \\
Anthropic & Claude Code, desarrolladores empresariales, capacidades Agentic & Si puede extenderse del coding a flujos de trabajo más amplios. \\
Thinking Machines & Multimodalidad nativa, visión de la interacción de nueva generación & Si puede convertir su línea de investigación en un producto sólido. \\
Meta/xAI & Plataformas, código abierto, activos de datos sociales o en tiempo real & Si puede crear un producto original de 0 a 1 que se gane un lugar en la mente de los usuarios. \\
\bottomrule
\end{tabularx}
\end{knowledgebox}

\subsection{Lo que la diferenciación enseña a los emprendedores}

Después del mapa de empresas, este apartado lo traslada a la perspectiva del emprendedor. Si las principales empresas empiezan a diferenciar sus rutas, el emprendedor ya no puede preguntarse solo «quién tiene el modelo más fuerte», sino «en qué cadena de capacidades me ubico». Quien construye herramientas de coding debe observar a Anthropic/Claude Code; quien trabaja en búsqueda e investigación, a Deep Research; quien hace interacción multimodal, a Thinking Machines y Gemini; quien apunta al punto de entrada del consumidor, al conjunto completo de productos de ChatGPT.

\begin{warningbox}{No interpretes la diferenciación como un aumento automático de oportunidades para las empresas pequeñas}
La diferenciación significa que las empresas líderes empiezan a elegir activamente sus campos de batalla, no que los abandonen. La oportunidad del emprendedor reside en ser más rápido, más específico y más profundo, no en enfrentarse de frente en los terrenos donde los gigantes ya entraron de forma explícita.
\end{warningbox}

\subsection{Resumen de la sección}

La diferenciación de las empresas de modelos indica que la industria entró en una etapa estratégica. Las rutas generalista, Coding+Agentic, multimodal y de ecosistema pueden producir ganadores, pero cada una exige capacidades distintas de producto, de datos y comerciales.
\section{Suite horizontal integral e integración vertical}

La sección anterior trató la diferenciación; esta trata la convergencia. En el mercado de consumo hay una tendencia clara a que todo converja en unos pocos líderes, y ChatGPT podría absorber muchos productos. El ejemplo de suite horizontal integral es ChatGPT, que reúne Chat, búsqueda, Coding, Agent y Workspace en una misma entrada; el ejemplo de integración vertical es Gemini, que, desde los chips TPU, el modelo Gemini y las aplicaciones Agent hasta Google Docs, Chrome, Android y YouTube, puede formar una superintegración.

\begin{figure}[H]
\centering
\includegraphics[width=0.94\textwidth]{horizontal-vs-vertical.png}
\caption{Suite horizontal integral frente a integración vertical: ChatGPT como suite horizontal, Gemini como superintegración vertical. Diagrama conceptual propio, elaborado a partir de la conversación de 00:21:37--00:33:35.}
\end{figure}

\begin{knowledgebox}{Lectura de la figura: lo horizontal y lo vertical son dos formas de absorber}
La suite horizontal integral concentra muchas tareas del usuario en una sola entrada; la integración vertical encadena chips, modelos, aplicaciones, navegador, sistema operativo y ecosistema de contenidos. La primera se apropia de la mente del usuario; la segunda, de la infraestructura y de las entradas predeterminadas.
\end{knowledgebox}

\subsection{La barrera de producto horizontal de ChatGPT}

ChatGPT no es solo un envoltorio de un modelo: tiene hábitos de los usuarios, posicionamiento de marca, una entrada predeterminada, historial de conversaciones y una expansión continua de funciones. Para un producto de IA independiente, lo más difícil es esto: si la función que construiste se convierte en una pestaña más de la suite de ChatGPT, ¿por qué el usuario seguiría usándote por separado?

\begin{importantbox}{El poder destructivo de la suite horizontal integral}
Lo más temible de la suite horizontal integral no es que cada función sea la mejor, sino que se convierte en la entrada predeterminada del usuario. Una entrada predeterminada puede distribuir con rapidez funciones «suficientemente buenas» a una enorme cantidad de usuarios y así acortar la ventana de oportunidad de los productos de las empresas emergentes.
\end{importantbox}

\subsection{La integración vertical de Gemini}

La ventaja de Google está en la integración vertical: TPU, nube, Gemini, búsqueda, publicidad, Docs, Chrome, Android y YouTube pueden funcionar de forma coordinada. Por eso mismo cambió la opinión de Guangmi (广密) sobre Google: cuando los productos de IA terminen por requerir la coordinación de búsqueda, publicidad, ofimática, navegador y sistema operativo, los activos antiguos de Google podrían volver a convertirse en ventajas nuevas.

\begin{knowledgebox}{Activos de la integración vertical de Google}
\begin{tabularx}{\textwidth}{p{0.22\textwidth}p{0.34\textwidth}X}
\toprule
Activo & Función & Valor en la era de la IA \\
\midrule
TPU/Cloud & Infraestructura de cómputo y de entrenamiento & Reduce el costo de los modelos y garantiza el suministro. \\
Search/Ads & Sistema de entrada y de monetización & La búsqueda con IA y las plataformas publicitarias podrían volver a fusionarse. \\
Chrome/Android & Navegador y sistema operativo móvil & Los Agent pueden entrar en el entorno digital predeterminado. \\
Docs/Workspace & Escenarios de oficina y archivos del usuario & Entrada a la productividad empresarial y personal. \\
YouTube & Contenido de video y datos/escenarios multimodales & Campo de aplicación para la comprensión y la generación multimodales. \\
\bottomrule
\end{tabularx}
\end{knowledgebox}

\subsection{Resumen de la sección}

La suite horizontal integral y la integración vertical explican por qué cada vez es más difícil esquivar a las empresas líderes. Los emprendedores deben convertirse en una capacidad clave dentro de la suite o encontrar escenarios que a los gigantes no les resulte fácil integrar.
\section{La inteligencia y el producto importan por igual}

Esta sección aborda la segunda palabra clave: producto. Durante los últimos tres años, la industria se volcó de forma casi obsesiva a explorar el límite superior de la inteligencia, pero ahora el producto vuelve a ser importante. ChatGPT tiene muchas barreras de entrada no técnicas, mientras que una herramienta de Coding o una empresa de modelos que solo cuenta con barreras técnicas puede ser alcanzada con facilidad. La fortaleza de OpenAI está en el equilibrio: por un lado explora el límite superior de la inteligencia y, por otro, convierte los beneficios de esa inteligencia en tráfico de producto y en presencia de marca en la mente de los usuarios.

\begin{figure}[H]
\centering
\includegraphics[width=0.96\textwidth]{intelligence-product-balance.png}
\caption{La inteligencia y el producto importan por igual: la exploración del límite superior de la inteligencia debe convertirse en tráfico, marca y presencia en la mente de los usuarios. Diagrama conceptual propio, elaborado a partir de la conversación de 00:33:35--00:38:52.}
\end{figure}

\begin{knowledgebox}{Lectura de la figura: la inteligencia no equivale automáticamente a un producto}
La capacidad del modelo aporta posibilidades; la experiencia de producto convierte esas posibilidades en uso predeterminado; el tráfico y la presencia de marca en la mente de los usuarios hacen que estos regresen de forma continua; la comercialización convierte la inteligencia en una capacidad de la empresa. Si falta cualquiera de estos eslabones, el resultado puede quedarse en una simple demo técnica.
\end{knowledgebox}

\subsection{Barreras del modelo y barreras del producto}

La sección anterior subrayó que el producto vuelve a ser importante; esta separa dos tipos de barreras. Las barreras del modelo provienen de los benchmark, el coding, las capacidades agentic, el contexto y el razonamiento; las barreras del producto provienen de los hábitos de los usuarios, la marca, los puntos de entrada, el retorno de datos y la comercialización. Las barreras del modelo pueden ser alcanzadas y las del producto también pueden ser absorbidas, pero solo la suma de ambas forma una empresa sólida.

\begin{figure}[H]
\centering
\includegraphics[width=0.94\textwidth]{model-vs-product-moat.png}
\caption{Barreras del modelo vs barreras del producto: las barreras técnicas pueden ser alcanzadas; las barreras del producto incluyen el tráfico, la marca y el punto de entrada predeterminado. Diagrama conceptual propio, elaborado a partir de la conversación de 00:33:35--00:38:52.}
\end{figure}

\begin{warningbox}{No hay que fijarse solo en las barreras técnicas}
Las barreras técnicas de un producto de IA pueden alcanzarse en un plazo de 3 a 6 meses. Si no hay retorno de datos, hábitos de los usuarios, canales, marca ni flujos de trabajo consolidados, la ventaja técnica puede convertirse fácilmente en una ventana pasajera.
\end{warningbox}

\begin{knowledgebox}{La parte no técnica de las barreras del producto}
\begin{tabularx}{\textwidth}{p{0.22\textwidth}p{0.34\textwidth}X}
\toprule
Barrera & Manifestación & Por qué es difícil de replicar \\
\midrule
Punto de entrada predeterminado & El usuario lo abre de forma natural todos los días & Requiere una marca y hábitos de uso construidos a largo plazo. \\
Integración en el flujo de trabajo & El producto entra en los procesos del equipo & El costo de sustitución es alto y los datos regresan de forma continua. \\
Presencia de marca en la mente del usuario & El usuario le asigna por omisión cierto tipo de tareas & El primer Magic Moment se recuerda durante mucho tiempo. \\
Comercialización & Cuenta con pagos y canales estables & No basta con un modelo potente; también hay que saber vender y presentar el producto. \\
\bottomrule
\end{tabularx}
\end{knowledgebox}

\subsection{Resumen de la sección}

La relación entre la inteligencia y el producto se ha vuelto más equilibrada. Las empresas de modelos necesitan la conversión en producto, y las empresas de producto también necesitan entender los modelos; ni la tecnología por sí sola ni el producto por sí solo son suficientes.
\section{Los productos de IA son como la minería: el Magic Moment y su ventana de vigencia}

Esta sección explica la metáfora de la «minería». El primer producto que ofrece una experiencia que asombra a los usuarios obtiene enormes beneficios de marketing; vale la pena aunque el consumo de token sea alto. Los primeros productos de preguntas y respuestas con búsqueda, Cursor y Manus son ejemplos de productos que aprovecharon un Magic Moment. Sin embargo, la ventana de oportunidad se está acortando: pasó de dos años a uno y luego a unos cuantos meses.

\begin{figure}[H]
\centering
\includegraphics[width=0.96\textwidth]{ai-product-mining-window.png}
\caption{Los productos de IA son como la minería: la ventana de vigencia del Magic Moment se está acortando. Diagrama conceptual propio, elaborado a partir de la conversación de 00:38:52--00:44:21.}
\end{figure}

\begin{knowledgebox}{Lectura de la figura: la minería no es una ventaja defensiva a largo plazo}
En la etapa de minería se compite por el primer asombro y por el beneficio de marca; cuando la ventana se acorta, las empresas de modelos, los competidores y las implementaciones de código abierto reaccionan con rapidez. Una ventaja defensiva real exige pasar del Magic Moment al flujo de trabajo, los datos, el punto de entrada y la marca.
\end{knowledgebox}

\subsection{¿Puede una empresa de producto ganarle a una empresa de modelos?}

La cuestión central de la ventana de minería es si una empresa de producto logra levantar barreras antes de que la empresa de modelos entre al mercado. La empresa de producto tiene ventajas de velocidad, enfoque y experiencia; la empresa de modelos tiene ventajas en modelos, distribución, capacidad de cómputo y un ecosistema completo de productos. Para ganar, la empresa de producto debe consolidar, dentro de la ventana, un posicionamiento sólido en la mente de los usuarios y un ciclo cerrado de datos; de lo contrario, la empresa de modelos la absorberá como una función más.

\begin{importantbox}{Fórmula de la ventana de oportunidad}
\[
\text{Startup chance} \approx \text{Magic Moment} \times \text{Speed} \times \text{Moat after copy}
\]
Donde Magic Moment es el primer asombro, Speed es la velocidad para aprovechar la ventana y Moat after copy es la barrera que se conserva después de que el producto ha sido copiado.
\end{importantbox}

\begin{knowledgebox}{Por qué se acorta la ventana}
La ventana de oportunidad se acorta por tres razones. Primero, las capacidades de los modelos base se actualizan más rápido, y la experiencia anterior pronto se vuelve una capacidad predeterminada. Segundo, las principales empresas de modelos tienen una distribución más fuerte y pueden copiar con rapidez los modelos de negocio exitosos. Tercero, el código abierto y las cadenas de herramientas aceleran la implementación técnica, de modo que el liderazgo de un producto no puede depender solo de una demo.
\end{knowledgebox}

\begin{knowledgebox}{Tabla de estrategias para emprendedores}
\begin{tabularx}{\textwidth}{p{0.22\textwidth}p{0.34\textwidth}X}
\toprule
Estrategia & Escenario de aplicación & Riesgo \\
\midrule
Aprovechar el Magic Moment & Acaba de surgir una nueva capacidad y los usuarios aún no tienen un punto de entrada predeterminado & La ventana es corta y es fácil que una empresa de modelos lo copie. \\
Especialización vertical profunda & Los flujos de trabajo del sector son complejos y los datos y las relaciones con clientes son difíciles de migrar & Crecimiento lento; ventas y entrega costosas. \\
Desarrollar cadenas de herramientas & Herramientas para desarrolladores o empresas construidas en torno a los modelos líderes & Los cambios en la API de la plataforma afectan al producto. \\
Convertirse en punto de entrada & Ganar los escenarios de uso predeterminados y el posicionamiento de marca & Competencia directa con ChatGPT/Gemini. \\
\bottomrule
\end{tabularx}
\end{knowledgebox}

\subsection{Resumen de la sección}

Los productos de IA son como la minería porque, cuando surge una nueva capacidad, se abre una breve ventana de altos rendimientos. Lo que debe hacer un emprendedor no es solo dar la primera palada, sino levantar barreras duraderas antes de que la ventana se cierre.
\section{Análisis de casos de la ventana de producto: del asombro a la herramienta predeterminada}

La sección anterior presentó la metáfora de la minería; esta sección la desglosa en el ciclo de vida de un producto para responder a la pregunta «¿por qué el Magic Moment deja de bastar tan pronto?». La trayectoria habitual de un producto de IA es la siguiente: surge una nueva capacidad de un modelo, un emprendedor la empaqueta en una primera experiencia que asombra, los usuarios la difunden con rapidez, la empresa del modelo y los competidores la replican y, al final, esa capacidad se convierte en una función predeterminada. Solo si, después de la difusión, el producto sigue acumulando flujos de trabajo, datos y marca, evita ser absorbido por la suite integral de la plataforma.

\begin{knowledgebox}{Ciclo de vida del Magic Moment}
\begin{tabularx}{\textwidth}{p{0.18\textwidth}p{0.34\textwidth}X}
\toprule
Etapa & Qué ocurre & Acción clave \\
\midrule
Descubrimiento de la capacidad & La nueva capacidad del modelo acaba de estar disponible & Encontrar pronto la tarea que mejor demuestre su valor. \\
Empaquetado de la experiencia & Se convierte en un Magic Moment que el usuario puede percibir & Reducir la barrera de uso y reforzar lo que impulsa la difusión. \\
Difusión explosiva & Se propaga por redes sociales y de boca en boca & Ocupar el lugar predeterminado en la mente del usuario y construir una marca temprana. \\
Reacción de la competencia & La empresa del modelo o productos similares la copian & Completar flujos de trabajo, datos y colaboración en equipo. \\
Conversión en función predeterminada & La capacidad pasa a ser una función básica de la plataforma & Migrar a escenarios más profundos o convertirse en un componente de la plataforma. \\
\bottomrule
\end{tabularx}
\end{knowledgebox}

\begin{importantbox}{La línea divisoria entre el asombro y la herramienta predeterminada}
El asombro nace de «ver por primera vez que puede hacerlo»; la herramienta predeterminada nace de «no poder prescindir de ella ningún día». Lo difícil para un producto de IA es completar esta transición antes de que se cierre la ventana del asombro.
\end{importantbox}

\begin{warningbox}{Efectos secundarios del Magic Moment}
El Magic Moment puede hacer creer al equipo que el producto ya está consolidado, pero el asombro del usuario no equivale a retención a largo plazo. El verdadero peligro es que el equipo siga optimizando en función del efecto de la demostración, en lugar de hacerlo en función de las tareas de alta frecuencia, los procesos de colaboración y las razones para pagar.
\end{warningbox}

\subsection{Resumen de la sección}

El ciclo de vida de los productos de IA es cada vez más corto, y los emprendedores deben llevar el producto de una experiencia que asombra a una herramienta predeterminada. De lo contrario, el Magic Moment no será más que una función dentro de la suite integral de la empresa del modelo.
\section{Experiencia de nivel L4: Deep Research y Claude Code}

Antes se habló de la ventana de minería; esta sección examina dónde está hoy la veta más clara. Guangmi (广密) considera que los dos mejores Agent ya ofrecen una experiencia L4: Deep Research de ChatGPT y Claude Code de Anthropic, que corresponden respectivamente a la búsqueda e investigación de información y al desarrollo de software. Hoy el mayor dividendo sigue siendo el de language/code, sobre todo el de code; todavía no es el de la multimodalidad, los modelos del mundo ni la robótica.

\begin{figure}[H]
\centering
\includegraphics[width=0.94\textwidth]{l4-experience-map.png}
\caption{Mapa de la experiencia L4: Deep Research y Claude Code representan la búsqueda y el desarrollo de software, respectivamente. Diagrama conceptual propio, elaborado a partir de la conversación de 00:44:21--00:52:43.}
\end{figure}

\begin{knowledgebox}{Lectura de la figura: por qué estos dos se parecen a L4}
Deep Research puede realizar búsquedas de información en varios pasos, sintetizarlas y redactar un informe; Claude Code puede entender un código base, modificarlo, probarlo e iterar. Todas estas tareas tienen una entrada clara, un entorno de herramientas y una salida verificable, por lo que es más fácil que den lugar a una experiencia de Agent de alto nivel.
\end{knowledgebox}

\subsection{Por qué el dividendo de code es el mayor}

El mundo del código tiene ventajas naturales: las tareas pueden formalizarse, las herramientas pueden invocarse, los resultados pueden probarse, la retroalimentación es rápida y el valor es alto. Frente a la multimodalidad, la robótica y los modelos del mundo, no sorprende que el coding haya alcanzado antes la experiencia L4. Es uno de los entornos de alto valor del mundo digital más adecuados para los Agent.

\begin{knowledgebox}{Condiciones de la experiencia L4}
\begin{tabularx}{\textwidth}{p{0.22\textwidth}p{0.34\textwidth}X}
\toprule
Condición & Deep Research & Claude Code \\
\midrule
Tarea clara & Preguntas de investigación y recopilación de fuentes & issue, bug, feature. \\
Herramientas disponibles & Búsqueda, páginas web, citas & Código base, pruebas, shell. \\
Resultado evaluable & Calidad del informe y de las fuentes & Pruebas aprobadas, review del código. \\
Valor alto & Ahorra tiempo de investigación & Ahorra tiempo de ingeniería. \\
\bottomrule
\end{tabularx}
\end{knowledgebox}

\begin{knowledgebox}{Próximos escenarios candidatos a L4}
\begin{tabularx}{\textwidth}{p{0.20\textwidth}p{0.34\textwidth}X}
\toprule
Escenario & Por qué podría surgir L4 & Dificultad \\
\midrule
Legal/cumplimiento & Uso intensivo de documentos, procesos claros, valor alto & Gran responsabilidad, alto costo de los errores. \\
Investigación financiera & Recuperación de información, modelos, informes y pistas de operaciones & Altas exigencias de licencias de datos y de tiempo real. \\
Operaciones de ventas & CRM, correo, reuniones y cotizaciones pueden convertirse en herramientas & Datos privados y procesos organizacionales complejos. \\
Análisis de datos & Ciclo cerrado de SQL, reportes, BI e interpretación del negocio & Requiere entender la definición de las métricas y la semántica del negocio. \\
Diseño/multimodal & La generación visual y la iteración de edición se aceleran & La evaluación subjetiva y los derechos de autor son más complejos. \\
\bottomrule
\end{tabularx}
\end{knowledgebox}

\subsection{Resumen de la sección}

La experiencia L4 no es un nivel abstracto, sino que el usuario esté realmente dispuesto a confiarle tareas complejas al sistema. Deep Research y Claude Code son hoy las dos muestras más claras.
\section{Reevaluar a Google: publicidad, búsqueda y la convergencia de su ecosistema de productos}

La sección anterior trató la experiencia L4 actual; esta sección vuelve a la competencia entre plataformas y, en particular, a la reevaluación de Google. Una hipótesis es que ChatGPT acabará construyendo una plataforma publicitaria, porque hace poco contrató a un nuevo CEO de comercialización; sin embargo, Google sigue siendo la mejor plataforma publicitaria del mundo. Al final, las formas de producto de todos podrían converger por caminos distintos: Search evoluciona, y la publicidad, la IA, el ecosistema de productos y el sistema operativo se fusionan poco a poco.

\begin{figure}[H]
\centering
\includegraphics[width=0.96\textwidth]{google-revaluation.png}
\caption{Reevaluar a Google: la lógica de la publicidad, la búsqueda, Gemini, las TPU y el ecosistema de productos vuelve a converger. Diagrama conceptual propio, elaborado a partir de la conversación entre 00:52:43 y 00:55:53.}
\end{figure}

\begin{knowledgebox}{Lectura de la figura: los activos antiguos de Google podrían convertirse en nuevas ventajas}
La búsqueda con IA necesita comercializarse, la publicidad necesita nuevos puntos de entrada, Gemini necesita distribución de producto y TPU/Cloud necesita demanda de modelos. Si Google recombina estos activos, obtendrá una ventaja de integración vertical.
\end{knowledgebox}

\begin{knowledgebox}{Comparación de estrategias de plataforma}
\begin{tabularx}{\textwidth}{p{0.22\textwidth}p{0.34\textwidth}X}
\toprule
Estrategia de plataforma & Enfoque central & Resultado posible \\
\midrule
OpenAI/ChatGPT & Integración horizontal de las tareas del usuario para convertirse en el punto de entrada predeterminado a la IA & Absorbe muchos de los usos ligeros de aplicaciones independientes. \\
Google/Gemini & Integración vertical de cómputo, búsqueda, publicidad, ofimática y OS & La IA reactiva los activos de la internet anterior. \\
Anthropic/Claude & Construir un posicionamiento técnico en la mente de los usuarios con la experiencia de Coding y Agentic & Liderazgo en escenarios de alto valor para desarrolladores y empresas. \\
Empresas emergentes verticales & Aprovechar la ventana de oportunidad en escenarios donde los gigantes no profundizan lo suficiente & Necesitan construir pronto barreras de datos y de flujos de trabajo. \\
\bottomrule
\end{tabularx}
\end{knowledgebox}

\subsection{Resumen de la sección}

El caso de Google nos recuerda que la IA no sustituirá sin más los activos de la internet anterior. La búsqueda, la publicidad, el navegador, la ofimática y el sistema operativo podrían ser integrados de nuevo por la IA.
\section{La burbuja, la AGI de los chinos y la narrativa global}

Las secciones anteriores analizan empresas y productos; la última sección trata las variables externas: el capital, la burbuja y la narrativa global. Si la AGI es una burbuja, qué podría hacerla estallar, en qué difiere la inteligencia humana de la de los gorilas, los nuevos temas en el Área de la Bahía y la frase «las finanzas, de los judíos; la AGI, de los chinos». Estos temas parecen dispersos, pero en realidad todos discuten la narrativa social de la IA y las expectativas del capital.

\begin{figure}[H]
\centering
\includegraphics[width=0.96\textwidth]{agi-bubble-dashboard.png}
\caption{Tablero de la AGI Bubble: que la burbuja estalle o no depende de los ingresos, el avance de los modelos, el CapEx y la materialización de los productos. Diagrama conceptual de elaboración propia, basado en la conversación de 00:55:53--01:09:11.}
\end{figure}

\begin{knowledgebox}{Lectura de la figura: la burbuja no se juzga solo por el ánimo del mercado}
Para evaluar la burbuja de la AGI hay que observar si los modelos siguen avanzando, si los ingresos se materializan, si el CapEx genera rendimiento, si siguen apareciendo productos L4 y si las valoraciones descuentan en exceso el futuro. El ánimo del mercado fluctúa, pero los fundamentos se juzgan con estos indicadores.
\end{knowledgebox}

\begin{knowledgebox}{Tabla de detonantes de la burbuja}
\begin{tabularx}{\textwidth}{p{0.22\textwidth}p{0.34\textwidth}X}
\toprule
Detonante & Posibles señales & Por qué es peligroso \\
\midrule
Desaceleración del avance de los modelos & Los modelos nuevos mejoran poco y cuestan más & Las valoraciones se basan en saltos continuos de capacidad. \\
Ingresos por debajo de lo esperado & Los ingresos por suscripciones, API y clientes empresariales son menores que el CapEx & El flujo de efectivo difícilmente justifica el gasto de capital. \\
Cierre de la ventana de producto & La experiencia L4 tarda en extenderse & Disminuyen el crecimiento de usuarios y la disposición a pagar. \\
Restricciones de hardware y energía & Suben los costos de GPU, centros de datos y electricidad & Los costos de entrenamiento e inferencia reducen los márgenes. \\
Riesgos regulatorios y de derechos de autor & Se restringen los datos y los resultados generados & Se ve afectado el suministro de datos y de productos. \\
\bottomrule
\end{tabularx}
\end{knowledgebox}

\subsection{La AGI de los chinos}

«Las finanzas, de los judíos; la AGI, de los chinos» es una síntesis con carga narrativa. Apunta a un hecho: en la competencia global por la AGI, los equipos de origen chino son muy activos en modelos, aplicaciones, ingeniería y emprendimiento. Pero los retos también son evidentes: la marca global, los mercados de capitales, el suministro de capacidad de cómputo, las ventas empresariales y el cumplimiento normativo no son problemas puramente técnicos.

\begin{figure}[H]
\centering
\includegraphics[width=0.94\textwidth]{chinese-agi-global.png}
\caption{La AGI de los chinos: en la competencia global por la AGI, se revalúan las capacidades técnicas, de producto y comerciales de los equipos de origen chino. Diagrama conceptual de elaboración propia, basado en la conversación de 00:55:53--01:09:11.}
\end{figure}

\begin{warningbox}{La narrativa no sustituye a la capacidad}
«La AGI de los chinos» puede ser una narrativa alentadora, pero no sustituye a la capacidad de los modelos, la capacidad de producto, la capacidad organizacional ni la capacidad comercial global. Una empresa verdaderamente global necesita resolver a la vez la tecnología, el producto, el capital, la marca y el cumplimiento normativo.
\end{warningbox}

\subsection{Resumen de la sección}

La burbuja y la narrativa global son variables externas del ciclo de la IA. El avance de los modelos y la experiencia de producto determinan la base; el capital y la narrativa determinan las fluctuaciones. Quien emprende debe observar ambos niveles a la vez.

\section{Síntesis y ampliación}

Esta sección condensa EP112 en cinco conclusiones. Primera: las empresas de modelos empiezan a diferenciarse, y lo general, Coding+Agentic, la interacción multimodal y el posicionamiento en el ecosistema se vuelven rutas distintas. Segunda: el punto de acceso para el consumidor converge, y la suite horizontal de ChatGPT y la integración vertical de Gemini representan dos formas de absorción. Tercera: tanto la inteligencia como el producto importan, y las barreras de producto empiezan a valorarse de nuevo. Cuarta: los productos de IA se parecen a la minería, y la ventana del Magic Moment es cada vez más corta. Quinta: la experiencia L4 ya apareció en Deep Research y Claude Code, y el mayor beneficio actual sigue estando en language/code.

\begin{importantbox}{EP112 dentro de la serie de IA de Zhang Xiaojun (张小珺)}
EP127 y EP136 son informes trimestrales más amplios sobre los modelos grandes; EP115/EP113 tratan sobre Agent y la ruta técnica de K2. EP112, en cambio, regresa la mirada al producto y al mercado: cómo se diferencian las empresas de modelos, cómo los productos de IA compiten por su ventana y cómo convergen las rutas de plataforma de Google/OpenAI/Anthropic.
\end{importantbox}

\subsection{Takeaways clave}

Las secciones anteriores abordaron por separado la diferenciación de las empresas, la convergencia de las plataformas, la ventana de producto y el relato de la burbuja; esta sección los condensa en juicios aplicables. Para el inversionista, lo importante no es seguir cada lanzamiento de un modelo nuevo, sino determinar si la capacidad se convierte en punto de acceso e ingresos; para el emprendedor, lo importante no es demostrar que el modelo puede hacerlo, sino demostrar que puede retener a los usuarios una vez cerrada la ventana.

\begin{knowledgebox}{Lista de juicios de este episodio}
\begin{tabularx}{\textwidth}{p{0.24\textwidth}p{0.34\textwidth}X}
\toprule
Juicio & Significado & Indicadores a observar \\
\midrule
Diferenciación de las empresas de modelos & Cada empresa elige una pila de capacidades distinta & Rankings de coding, agentic, multimodalidad y uso general. \\
El producto vuelve a importar & La inteligencia debe convertirse en punto de acceso y en posicionamiento en la mente del usuario & DAU, retención, pago, elección predeterminada de marca. \\
La ventana se acorta & El Magic Moment se copia más rápido & Velocidad de respuesta de la competencia, absorción en funciones nativas del modelo. \\
Google se revalúa & Activos antiguos pueden recombinarse en la IA & Integración de Search, Ads, Gemini y Workspace. \\
La burbuja depende de que se cumpla lo prometido & El ánimo del mercado no es el único criterio & Ingresos, CapEx, avance de los modelos y productos L4. \\
\bottomrule
\end{tabularx}
\end{knowledgebox}

\begin{knowledgebox}{Marco de decisión para invertir/emprender}
\begin{tabularx}{\textwidth}{p{0.22\textwidth}p{0.34\textwidth}X}
\toprule
Pregunta & Qué observa el inversionista & Qué debe responder el emprendedor \\
\midrule
¿La capacidad es escasa? & ¿Cuánto tardarán las empresas de modelos en alcanzarla? & ¿Cuánto dura mi ventana de experiencia? \\
¿El punto de acceso es claro? & ¿Los usuarios lo abren con frecuencia? & ¿Puedo convertirme en el flujo de trabajo predeterminado? \\
¿Los datos regresan al producto? & ¿Mejora mientras más se usa? & ¿Puedo acumular datos o preferencias propios? \\
¿El modelo de negocio se sostiene? & ¿Los ingresos cubren el costo de inferencia? & ¿Cómo convierto el Magic Moment en una razón para pagar? \\
¿Lo harán los gigantes? & ¿La suite completa ya lo incluye de forma natural? & ¿Puedo evitar la plataforma o apoyarme en ella? \\
\bottomrule
\end{tabularx}
\end{knowledgebox}

\begin{enumerate}
\item La diferenciación de las empresas de modelos no es un debilitamiento, sino una elección de ruta: lo general, Coding+Agentic y la interacción multimodal compiten por posiciones dominantes distintas.
\item La suite horizontal y la integración vertical seguirán reduciendo el espacio de supervivencia de los productos de IA independientes.
\item Las barreras de producto incluyen el punto de acceso predeterminado, el posicionamiento de marca, el retorno de datos y la monetización, no solo la capacidad del modelo.
\item La Magic Window de los productos de IA se está acortando; el emprendedor debe convertir rápidamente el asombro en una barrera.
\item La experiencia L4 más clara sigue estando en language/code, no en la multimodalidad, los modelos del mundo ni la robótica.
\end{enumerate}

\subsection{Preguntas abiertas}

\begin{enumerate}
\item ¿La próxima experiencia L4 aparecerá en derecho, finanzas, ventas, análisis de datos o en el workflow empresarial?
\item ¿Puede la integración vertical de Google contrarrestar la ventaja de punto de acceso horizontal de ChatGPT?
\item ¿Cómo puede una empresa de producto construir una barrera duradera dentro de una Magic Window de 3 a 6 meses?
\item Si se forma una burbuja de la AGI, ¿el punto de ruptura más probable será la desaceleración del avance de los modelos, ingresos por debajo de lo esperado o un rendimiento insuficiente del CapEx?
\item Una vez que las empresas de modelos conviertan sus capacidades en productos, ¿en qué capas conviene más que permanezcan las empresas de aplicaciones independientes: el punto de acceso, la cadena de herramientas, el flujo de trabajo vertical o la red de datos?
\item Cuando se fusionen la búsqueda, la publicidad, la ofimática y Agent, ¿la competencia entre Google y OpenAI se parecerá más a una guerra de buscadores o a una guerra de sistemas operativos?
\end{enumerate}

\begin{warningbox}{Vigencia de los juicios trimestrales}
Los juicios trimestrales caducan por naturaleza. Lo que de verdad vale la pena conservar no es el liderazgo momentáneo de una empresa, sino el marco de análisis: diferenciación, convergencia, ventana de producto, experiencia L4, integración de plataformas e indicadores de burbuja. Este marco puede seguir usándose para observar el próximo trimestre.
\end{warningbox}

\subsection{Lecturas adicionales}

\begin{itemize}
\item Si le interesa la ruta técnica de Agent, puede contrastarse con la entrevista a Yao Shunyu (姚顺雨) en EP115 y la entrevista a Yang Zhilin (杨植麟) sobre K2 en EP113.
\item Si le interesa la adopción en empresas, puede contrastarse con la entrevista a Wu Minghui (吴明辉) en EP116, para entender qué forma toma el Agentic Model con datos privados en el ámbito ToB.
\item Si le interesan los informes trimestrales más amplios sobre los modelos grandes, puede continuar con las notas de los informes trimestrales de EP127 y EP136.
\end{itemize}

\end{document}
